\section{DAO Key Migration}
\label{sec:dao-migration}

The Zoo \DAO{} authorises treasury disbursements, governance votes,
contract upgrades, and per-\LLM{} chain admissions. In 2.x the
\DAO{}'s authority surface used classical multisig and ECDSA. In 3.0
the entire \DAO{} authority moves to threshold post-quantum signing.

\subsection{What moves}

\begin{itemize}[leftmargin=1.2em]
  \item \textbf{Treasury custody key.} Threshold $\MLDSA{}\text{-}65$
        on M-Chain (per LP-017 schedule \cite{lp-017}).
  \item \textbf{Governance proposal-signing key.} Threshold
        \Ringtail{} for proposal hashes (PQ-safe with smaller
        per-signer state than $\MLDSA{}$, useful for high-frequency
        proposal flows).
  \item \textbf{Contract-upgrade authority key.} Threshold
        $\MLDSA{}\text{-}65$ with optional human-in-the-loop second
        factor (a hardware-key-attested \MLDSA{} from each
        upgrade-authorised member).
  \item \textbf{Per-\LLM{} chain admission key.} Per-chain
        threshold $\MLDSA{}$, scoped so that a compromise of one
        chain's admission key does not propagate to others.
\end{itemize}

\subsection{Why threshold ML-DSA + Ringtail and not just one}

The \DAO{} runs two independent PQ schemes for the same reason as
the \Quasar{} 3.0 cert design (\S\ref{sec:triple-cert}): two
different lattice problems, so a structural break in one does not
unlock the other. Treasury custody --- the highest-stakes authority
--- uses $\MLDSA{}\text{-}65$ because FIPS 204 conformance maximises
ecosystem auditability. Proposal-signing uses \Ringtail{} because
its smaller per-signer state suits high-frequency vote production.

\subsection{Ceremony for the migration}

The \DAO{} cannot simply self-rotate to PQ keys without a public
ceremony, because the legacy ECDSA holders need an auditable transfer
event. The ceremony:

\begin{enumerate}[leftmargin=1.4em]
  \item \textbf{Pre-announcement.} A governance proposal, signed
        under the legacy multisig, announces the migration window
        and target threshold parameters.
  \item \textbf{Distributed key generation (DKG).} Each member runs
        a $\MLDSA{}$ DKG to produce the new threshold key shares.
        Output is the threshold public key plus per-member secret
        shares.
  \item \textbf{Co-signed handover.} A single transaction signed by
        \emph{both} the legacy ECDSA multisig and the new threshold
        $\MLDSA{}$ key transfers custody from the legacy address
        to the threshold address.
  \item \textbf{Sunset.} The legacy multisig is retired after a
        defined window during which any member can challenge the
        ceremony with a fraud proof.
\end{enumerate}

A failed ceremony reverts to the legacy multisig and the migration
proposal must be re-issued. The ceremony is itself anchored to
\Quasar{} 3.0 cert \cite{lp-020}, providing triple-cert finality on
the handover record.

\subsection{Vote-signing migration}

Voter wallets in 3.0 carry dual keys: a legacy ECDSA key for backwards
compatibility, and a new $\MLDSA{}\text{-}65$ key. During the sunset
period both signature types count toward quorum. After sunset, only
$\MLDSA{}$ ballots count; legacy-only voters who have not co-signed
a wallet upgrade are silently dropped from the eligible set.

The wallet's UX presents the dual-key state explicitly so users
understand which signature scheme is in use for each action.

\subsection{Optional FHE ballots}

Per the holographic-consensus \DAO{} design \cite{zoo-2025-securities-and-dao},
ballots can be cast in \FHE{}-encrypted form so vote direction is
private until tally. In 3.0 the encryption is layered with the PQ
signature: the ballot is FHE-encrypted, then the FHE ciphertext is
signed under the voter's $\MLDSA{}$ key. The tally process operates
over ciphertexts; the result is decrypted only at the published tally
round. Voter authentication remains PQ-safe even though vote content
is privacy-preserving.

\subsection{Recovery in case of share loss}

Threshold key share loss is a known operational hazard. The 3.0
recovery procedure:

\begin{itemize}[leftmargin=1.2em]
  \item If $\leq n - t$ shares are lost, the threshold remains
        operable. Lost members are removed via governance and the
        committee re-keys at the next rotation epoch.
  \item If more than $n - t$ shares are lost, the threshold is
        broken; recovery requires a governance-authorised
        re-DKG signed by a quorum of the remaining live members
        plus a delay-locked failsafe authority that can authorise
        re-key after a multi-week timelock if even the remaining
        live members are insufficient.
  \item All recovery actions are publicly logged and anchored to L1
        cert.
\end{itemize}

\subsection{Effect on per-LLM chain governance}

Per-\LLM{} chains \cite{zoo-2025-securities-and-dao} run their own
local governance for model-specific decisions (parameter votes,
experience-library admissions, dataset whitelist changes).
Each per-\LLM{} chain inherits the 3.0 governance template: dual-key
ballots, threshold $\MLDSA{}$ for binding decisions, optional \FHE{}
voting. This means a chain compromise on one \LLM{} chain does not
require a global Zoo \DAO{} re-key.
